\documentclass[11pt,a4paper]{article}

\usepackage[margin=2.5cm]{geometry}
\usepackage{amsmath,amssymb,physics}
\usepackage{graphicx}
\usepackage{booktabs}
\usepackage{hyperref}
\usepackage{xcolor}
\usepackage{caption}
\usepackage{subcaption}
\usepackage{siunitx}
\usepackage{float}

\hypersetup{colorlinks=true,linkcolor=blue,citecolor=blue,urlcolor=blue}

\newcommand{\Fproj}{F_{\rm proj}}
\newcommand{\Lavg}{L_{\rm avg}}
\newcommand{\Utarget}{U_{\rm target}}
\newcommand{\SQR}{\mathrm{SQR}}
\newcommand{\SNAP}{\mathrm{SNAP}}
\newcommand{\CP}{\mathrm{CP}}
\newcommand{\disp}{\mathrm{D}}

\title{\textbf{cQED Synthesis of a Hybrid Qubit--Cavity Unitary}\\[6pt]
\large Numerical Study: Phases 1--6}
\author{Shyam Shankar Quantum Circuits Group}
\date{\today}

\begin{document}
\maketitle
\tableofcontents
\newpage

%============================================================
\section{Introduction and Target Gate}
%============================================================

We study the synthesis of a two-qubit-equivalent entangling unitary on a
dispersive circuit-QED platform.  The target gate, defined on the logical
subspace $\mathcal{H}_L = \mathrm{span}\{\ket{g,0},\ket{g,1},\ket{e,0},\ket{e,1}\}$
with qubit-first ordering (flat index $q N + n$, with $N$ the cavity truncation
dimension), is
\begin{equation}
  \Utarget = \frac{1}{\sqrt{2}}
  \begin{pmatrix} 1 & 0 & 1 & 0 \\ 1 & 0 & -1 & 0 \\
                  0 & 1 & 0 & 1 \\ 0 & -1 & 0 & 1 \end{pmatrix},
  \label{eq:utarget}
\end{equation}
expressed in the ordered basis $\{\ket{g,0},\ket{g,1},\ket{e,0},\ket{e,1}\}$.

\paragraph{Device parameters.}
\begin{align*}
  \chi  &= 2\pi\times(-2.84~\text{MHz}), &
  \chi' &= 2\pi\times(-21~\text{kHz}),  &
  K     &= 2\pi\times(-28~\text{kHz}),
\end{align*}
with dispersive Hamiltonian
$H/\hbar = -(\chi/2)\hat{a}^\dagger\hat{a}\,\sigma_z
           - (\chi'/2)(\hat{a}^\dagger\hat{a})^2\,\sigma_z
           + (K/2)(\hat{a}^\dagger\hat{a})^2$,
where $\sigma_z|g\rangle=+|g\rangle$, $\sigma_z|e\rangle=-|e\rangle$.
The characteristic time scales are $t_\pi = \pi/|\chi| = 0.176$~\si{\micro s},
$\chi'/\chi = 0.74\%$, and $K/\chi = 0.99\%$.

\paragraph{Primitive libraries.}
Two gate libraries are studied:
\begin{itemize}
  \item \textbf{Library 1}: $\{\disp(\alpha),\,R_q(\theta,\phi),\,\SQR_n(\theta_n,\phi_n)\}$
        — Displacement provides Fock-level mixing; QubitRotation $R_q$ gives
        unconditional qubit rotations; $\SQR_n$ gives Fock-conditioned qubit rotations.
  \item \textbf{Library 2}: $\{\disp(\alpha),\,R_q(\theta,\phi),\,\text{ConditionalPhaseSQR}\}$
        — Replaces the full SQR rotation with a conditional phase gate that applies
        Fock-number-dependent phases without rotating the qubit.
\end{itemize}

%============================================================
\section{Phase 1: Algebraic Structure of \texorpdfstring{$\Utarget$}{U\_target}}
%============================================================

\subsection{Unitarity and Gate Classification}

All 38 algebraic checks passed:
\begin{itemize}
  \item $\|U^\dagger U - I\|_F = 4.5\times10^{-16}$, $\det(U)=+1$.
  \item Minimum Schmidt singular value $= 1.0$: $\Utarget$ is \emph{maximally entangling}.
\end{itemize}

\subsection{CNOT Factorization}

$\Utarget$ admits the exact three-gate factorization
\begin{equation}
  \Utarget = (I_q \otimes H_c)\cdot \mathrm{CNOT}_{c\to q}\cdot \mathrm{CNOT}_{q\to c},
  \label{eq:factorization}
\end{equation}
where $H_c$ is the logical cavity Hadamard (swapping $\ket{n=0}\leftrightarrow\ket{n=1}$),
$\mathrm{CNOT}_{c\to q}$ flips the qubit conditioned on the cavity being in $\ket{n=1}$,
and $\mathrm{CNOT}_{q\to c}$ flips the cavity conditioned on the qubit being in $\ket{e}$.
The decomposition error is $\leq 10^{-15}$ (machine precision).

\paragraph{SQR nativity.}
\begin{itemize}
  \item $\mathrm{CNOT}_{c\to q}$: \textbf{native} — $\SQR_1(\pi,0)$ implements
        $-iX$ conditioned on $n=1$; the CP correction
        $\text{CP}([0,\pi/2])$ corrects the $-i$ phase to give exact CNOT,
        verified at $\|U-\mathrm{CNOT}_{c\to q}\|_F = 0$.
  \item $\mathrm{CNOT}_{q\to c}$ (qubit-controlled cavity level flip):
        \textbf{not SQR-native} — requires Fock-level mixing, which SQR gates cannot
        provide.  Displacement ($\disp$) is the required primitive.
  \item $H_c$ (cavity logical Hadamard): also requires Fock mixing; not SQR-native.
\end{itemize}

\textbf{Consequence:} Any gate library that omits Displacement cannot synthesize
$\Utarget$; its maximum achievable fidelity with SQR+CP alone is $F=0.5$
(verified numerically in Phase~2).

\begin{figure}[H]
  \centering
  \includegraphics[width=0.7\linewidth]{../figures/phase1_subgates.png}
  \caption{Phase 1: Verification of the three sub-gates ($\mathrm{CNOT}_{c\to q}$,
  $\mathrm{CNOT}_{q\to c}$, $H_c$) and the full CNOT factorization at
  $N=2$ and $N=8$.  All checks pass with numerical error $\leq 10^{-15}$.}
  \label{fig:phase1}
\end{figure}

%============================================================
\section{Phase 2: Idealized Synthesis in the 4D Logical Space}
%============================================================

\subsection{Setup}

We run \texttt{UnitarySynthesizer} in the 4D space (cavity truncation $N=2$,
$\mathrm{FULL\_DIM}=4$) with no dispersive drift (instantaneous gates).
Subspace: \texttt{Subspace.custom(4, [0,1,2,3], labels)}.

\subsection{Key Results}

\begin{table}[H]
  \centering
  \caption{Phase 2 synthesis results (ideal, $N=2$, no drift).}
  \label{tab:phase2}
  \begin{tabular}{lccl}
    \toprule
    Test & $\Fproj$ & Objective & Notes \\
    \midrule
    A: Analytic 3-gate (exact PrimitiveGate) & 1.000\,000 & $2\times10^{-16}$
      & Validates simulator + subspace \\
    B: D+SQR+CP depth-9 (multistart=16) & 1.000\,000 & $2\times10^{-15}$
      & Optimizer achieves exact synthesis \\
    C: SQR+CP only (no Displacement)     & 0.500\,000 & 0.50
      & Physics limit; Fock mixing absent \\
    \midrule
    D: Depth sweep (D+SQR+CP library) & & & \\
    \quad depth = 1 & 0.7071 & & $=1/\sqrt{2}$, insufficient \\
    \quad depth = 2 & 1.000\,000 & & \textbf{minimum depth for $F=1$} \\
    \quad depth = 3, 4 & 1.000\,000 & & Redundant \\
    \bottomrule
  \end{tabular}
\end{table}

\paragraph{Physical interpretation.}
The SQR+CP plateau at $F=0.5$ is a fundamental quantum control limit: SQR gates
implement qubit rotations conditioned on cavity Fock number $n$ but cannot transfer
population between Fock levels.  Since $\Utarget$ requires $H_c$ (a Hadamard mixing
$\ket{0}\leftrightarrow\ket{1}$) and $\mathrm{CNOT}_{q\to c}$ (both Fock-level
mixing operations), Displacement is a \emph{necessary} primitive.

\begin{figure}[H]
  \centering
  \includegraphics[width=\linewidth]{../figures/phase2_synthesis.png}
  \caption{Phase 2: Comparison of ansatz families and the depth sweep showing
  minimum depth = 2 blocks of D+SQR+CP for unit fidelity synthesis.}
  \label{fig:phase2}
\end{figure}

%============================================================
\section{Phase 3: Embedding in Physical Cavity ($N=8$)}
%============================================================

\subsection{Truncation Convergence}

The analytic decomposition \eqref{eq:factorization} is embedded in the full
$2N$-dimensional Hilbert space for $N\in\{2,4,6,8,10,12\}$.

\begin{table}[H]
  \centering
  \caption{Phase 3 truncation convergence (analytic decomposition).}
  \label{tab:phase3_trunc}
  \begin{tabular}{ccccc}
    \toprule
    $N$ & $F$ & $\|U_{\rm log}-U_{\rm target}\|_F$
        & $L_{\rm SQR_1}$ & $L_{\rm CNOT_{q\to c}}$ \\
    \midrule
    2 & 1.000\,000 & 0 & 0 & 0 \\
    4 & 1.000\,000 & 0 & 0 & 0 \\
    6 & 1.000\,000 & 0 & 0 & 0 \\
    8 & 1.000\,000 & 0 & 0 & 0 \\
    10& 1.000\,000 & 0 & 0 & 0 \\
    12& 1.000\,000 & 0 & 0 & 0 \\
    \bottomrule
  \end{tabular}
\end{table}

The analytic gate sequence has zero leakage for all $N$ because the sub-gate
matrices are constructed directly in the full $2N$-dimensional space.

\subsection{Displacement Leakage}

Displacement $\disp(\alpha)$ coherently populates higher Fock levels:
\begin{table}[H]
  \centering
  \caption{Average leakage $L_{\rm avg}$ of a single Displacement gate.}
  \label{tab:disp_leak}
  \begin{tabular}{cc}
    \toprule
    $|\alpha|$ & $L_{\rm avg}$ \\
    \midrule
    0.1 & 0.010 \\
    0.3 & 0.076 \\
    0.5 & 0.159 \\
    1.0 & 0.308 \\
    \bottomrule
  \end{tabular}
\end{table}

\subsection{Higher-order Dispersive Drift}

The overlap between free evolution under the full Hamiltonian
($\chi,\chi',K$) and the $\chi$-only approximation remains $F=1.000$ for all
times up to $4t_\pi = 0.70$~\si{\micro s}.  The ratio $\chi'/\chi = 0.74\%$
and $K/\chi = 0.99\%$ are sufficiently small that higher-order dispersive terms
are negligible over the relevant gate timescales.

\begin{figure}[H]
  \centering
  \includegraphics[width=\linewidth]{../figures/phase3_embedding.png}
  \caption{Phase 3: (left) truncation convergence and displacement leakage;
  (right) dispersive drift fidelity (full vs.\ $\chi$-only) over time.}
  \label{fig:phase3}
\end{figure}

%============================================================
\section{Phase 4: Variational Synthesis with Physical Drift}
%============================================================

\subsection{Setup}

Cavity truncation $N=8$ ($\mathrm{FULL\_DIM}=16$).
Subspace: \texttt{Subspace.custom(16, [0,1,8,9], labels)}.
Full dispersive drift model $(\chi,\chi',K)$ active during SQR and CP gates.
Optimization hyperparameters: \texttt{multistart}=2, \texttt{maxiter}=150,
$\lambda_{\rm leak}=0.1$.
Each library is tested at three depths.  Depth refers to the total number of
gates in the sequence.

\subsection{Results}

\begin{table}[H]
  \centering
  \caption{Phase 4 variational synthesis results (physical drift, $N=8$).}
  \label{tab:phase4}
  \begin{tabular}{lcccc}
    \toprule
    Ansatz & Depth & $\Fproj$ & $\Lavg$ & $L_{\rm worst}$ \\
    \midrule
    \multicolumn{5}{l}{\textit{Library 1: $\{\disp,\,R_q,\,\SQR\}$}} \\
    L1a: D R SQR D R SQR D     & 7  & 0.8846 & 0.1226 & 0.1562 \\
    L1b: D R SQR R D R SQR R D & 9  & 0.8971 & 0.1225 & 0.1401 \\
    \textbf{L1c: depth-11}     & \textbf{11} & \textbf{0.9193} & \textbf{0.1014} & \textbf{0.1364} \\
    \midrule
    \multicolumn{5}{l}{\textit{Library 2: $\{\disp,\,R_q,\,\CP\}$}} \\
    L2a: D R CP D R CP D       & 7  & 0.6726 & 0.0921 & 0.1168 \\
    L2b: D R CP R D R CP R D   & 9  & 0.6858 & 0.3523 & 0.4613 \\
    L2c: depth-11              & 11 & 0.8741 & 0.1236 & 0.1607 \\
    \bottomrule
  \end{tabular}
\end{table}

\paragraph{Best result.}
Library 1 (D+$R_q$+SQR) achieves the highest fidelity at all depths,
with the depth-11 ansatz L1c reaching $\Fproj=0.9193$ and $\Lavg=0.101$.
Fidelity increases monotonically with depth for both libraries, indicating
the optimizer has not yet saturated the expressiveness of either ansatz family.

\paragraph{Library 2 comparison.}
Library 2 (D+$R_q$+ConditionalPhaseSQR, abbreviated CP) underperforms Library 1
at every depth.  The CP gate applies only Fock-number-selective phases to the
qubit state without performing a full Bloch-sphere rotation, so it provides
fewer degrees of freedom for compensation of the dispersive drift pattern.
The depth-9 Library 2 ansatz exhibits anomalously high leakage ($\Lavg=0.352$),
indicating that the optimizer converges to a local minimum in that ansatz
topology.  At depth-11 the gap narrows ($\Fproj=0.874$ vs.\ $0.919$), but
Library 1 remains superior.

\paragraph{Effect of dispersive drift.}
The SQR gates accumulate dispersive phase
$\phi_n = \chi n t_{\rm SQR} + \chi' n(n-1) t_{\rm SQR} + K n(n-1) t_{\rm SQR}/2$
during their 2-\si{\micro s} execution.  The variational optimizer partially
compensates this drift by tuning the SQR azimuthal angles $\phi_n$; the
residual infidelity comes from leakage into $n\geq 2$ Fock states (see
Phase~5).

\begin{figure}[H]
  \centering
  \includegraphics[width=\linewidth]{../figures/phase4_synthesis_results.png}
  \caption{Phase 4: fidelity and leakage for all six ansatz variants
  (blue = Library 1 D+$R_q$+SQR, orange = Library 2 D+$R_q$+CP).}
  \label{fig:phase4}
\end{figure}

%============================================================
\section{Phase 5: Finite-Duration Pulse Effects}
%============================================================

\subsection{Drift During SQR Gates}

We compare the optimized L1c sequence (best from Phase~4) evaluated with
(a) the full drift model and (b) no drift during SQR gates (instantaneous
gate approximation):

\begin{table}[H]
  \centering
  \caption{Phase 5: physical vs.\ ideal drift during gates (L1c parameters).}
  \label{tab:phase5_drift}
  \begin{tabular}{lcc}
    \toprule
    Model & $\Fproj$ & $\Lavg$ \\
    \midrule
    Physical (drift during 2-\si{\micro s} SQR) & 0.91926 & 0.1014 \\
    Ideal (instantaneous gates)              & 0.91926 & 0.1014 \\
    \midrule
    Infidelity penalty from drift & 0.000\,00 & \\
    \bottomrule
  \end{tabular}
\end{table}

\paragraph{Interpretation.}
The dispersive phase accumulated during SQR execution,
$\phi_{n} = \chi n t + \chi' n(n-1) t + K n(n-1) t/2$,
is \emph{completely absorbed} by the variational optimizer through tuning of
the SQR azimuthal angles $\phi_n$.  The residual
infidelity $1-F = 8.07\%$ originates entirely from population leakage into
$n\geq 2$ Fock states (quantified by $\Lavg=10.14\%$), not from drift
compensation failure.  The higher leakage relative to SNAP-based sequences
(Phase~4, prior results $\Lavg=7.3\%$) reflects the absence of independent
Fock-level phase control in the Library 1 ansatz.

\subsection{SQR Duration and Higher-Order Phase Accumulation}

For Fock level $n=2$, the excess dispersive phase beyond the linear-in-$n$
term (which is absorbed by $\phi_n$ tuning) grows as
\begin{equation}
  \delta\phi_{n=2} = (2\chi' + K)\,t_{\rm SQR},
\end{equation}
giving $|\delta\phi_{n=2}| = 0.88$~rad at $t_{\rm SQR}=2$~\si{\micro s}.
This directly drives $n=2$ leakage and is the dominant source of residual
infidelity.

\subsection{Coherence Budget}

The L1c sequence total duration is
$T_{\rm seq} = 3\times200\,\text{ns} + 5\times100\,\text{ns} + 3\times2000\,\text{ns}
= 7.1$~\si{\micro s}.
\begin{table}[H]
  \centering
  \caption{Phase 5: coherence budget for L1c sequence.}
  \label{tab:phase5_budget}
  \begin{tabular}{lcc}
    \toprule
    Coherence channel & Typical time & Loss at $T_{\rm seq}$ \\
    \midrule
    Qubit $T_1$  & 9.8~\si{\micro s} & 51.5\% \\
    Qubit $T_2$  & 6.3~\si{\micro s} & 67.6\% \\
    Cavity $T_1$ & 200~\si{\micro s} & 3.5\%  \\
    \bottomrule
  \end{tabular}
\end{table}

The 7.1-\si{\micro s} sequence duration (vs.\ 9.6~\si{\micro s} for a
SNAP-based 9-gate sequence) provides a modest coherence advantage.
Decoherence remains the primary experimental challenge, since the gate time
is still comparable to the qubit $T_1$.

\begin{figure}[H]
  \centering
  \includegraphics[width=\linewidth]{../figures/phase5_pulse_effects.png}
  \caption{Phase 5: (left) ideal vs.\ physical fidelity; (center) dispersive
  phase at $n=2$ vs.\ SQR duration; (right) gate duration breakdown with $T_1$
  reference.}
  \label{fig:phase5}
\end{figure}

%============================================================
\section{Phase 6: Robustness Analysis}
%============================================================

\subsection{Setup}

We perturb the optimized L1c parameters about their nominal values and
re-evaluate $\Fproj$ without re-optimizing.  Perturbations cover $\chi$
miscalibration, displacement/SQR amplitude scaling, pulse duration scaling,
and global phase offset.  The nominal fidelity is $F_0 = 0.9193$.

\begin{table}[H]
  \centering
  \caption{Phase 6: robustness of L1c to calibration errors.}
  \label{tab:phase6}
  \begin{tabular}{lccccc}
    \toprule
    Perturbation & Range & $F$ at $-5\%$ & $F$ at $0$ & $F$ at $+5\%$
                 & RMS drop (±5\%) \\
    \midrule
    $\chi$ error           & ±10\% & 0.91926 & 0.91926 & 0.91926 & 0.0000 \\
    Amplitude error        & ±10\% & 0.91075 & 0.91926 & 0.90664 & 0.0058 \\
    Duration error         & ±15\% & 0.91926 & 0.91926 & 0.91926 & 0.0000 \\
    Phase offset $\delta\phi$ & ±0.2 rad
      & $F(-0.1)=0.91705$ & 0.91926 & $F(+0.1)=0.91688$ & 0.0012 \\
    \bottomrule
  \end{tabular}
\end{table}

\paragraph{Key findings.}
\begin{enumerate}
  \item \textbf{$\chi$ calibration}: zero fidelity sensitivity.  The dispersive
        drift phase in the logical $\{n=0,n=1\}$ subspace is fully absorbed
        by the SQR azimuthal angles $\phi_n$, making L1c intrinsically robust
        to $\chi$ miscalibration (same mechanism as Phase~5).
  \item \textbf{Amplitude}: RMS drop 0.0058 over ±5\% range.  Very robust.
  \item \textbf{Duration}: zero sensitivity, by the same $\chi t$ scaling argument.
  \item \textbf{Phase offset}: RMS drop 0.0012 over ±0.1~rad range.  Excellent
        tolerance.
\end{enumerate}

\begin{figure}[H]
  \centering
  \includegraphics[width=\linewidth]{../figures/phase6_robustness.png}
  \caption{Phase 6: robustness sweeps for the L1c optimised sequence.
  The horizontal dashed line marks $F_0=0.9193$.}
  \label{fig:phase6}
\end{figure}

%============================================================
\section{Summary and Conclusions}
%============================================================

\subsection{Main Results}

\begin{table}[H]
  \centering
  \caption{Summary of main findings across all phases.}
  \label{tab:summary}
  \begin{tabular}{lll}
    \toprule
    Phase & Key Result & Status \\
    \midrule
    1: Algebraic & $\Utarget$ is maximally entangling;
      CNOT factorization exact & \textcolor{green!60!black}{PASS (38/38)} \\
    2: Ideal 4D  & D+SQR+CP achieves $F=1$ at depth-2;
      SQR+CP limited to $F=0.5$ & Confirmed \\
    3: Embedding & Truncation convergence at $N=6$;
      analytic $F=1$ for all $N$ & Confirmed \\
    4: Physical  & L1c (D+$R_q$+SQR, depth-11) best: $\Fproj=0.9193$, $\Lavg=0.101$
      & Main result \\
    5: Pulses    & Drift fully compensated; residual
      infidelity = leakage only & Explained \\
    6: Robustness & Zero $\chi$/duration sensitivity;
      amplitude RMS 0.006, phase RMS 0.001 & Robust \\
    \bottomrule
  \end{tabular}
\end{table}

\subsection{Synthesis Feasibility}

$\Utarget$ is synthesizable in a dispersive cQED system:
\begin{itemize}
  \item \textbf{Ideally}: $F=1$ is achievable with depth-2 D+SQR+CP sequences
        in the 4D logical subspace (Phase~2).
  \item \textbf{Physically}: with an 11-gate D+$R_q$+SQR sequence (Library 1)
        and physical drift $(\chi,\chi',K)$, the best optimizer-found solution
        achieves $\Fproj=0.9193$ with $\Lavg=10.1\%$ at $N=8$.
  \item \textbf{Library hierarchy}: Library 1 (D+$R_q$+SQR) outperforms
        Library 2 (D+$R_q$+CP) at all depths because the SQR gate's full
        Bloch-sphere rotation provides more expressive qubit control than
        the CP phase-only gate.
  \item \textbf{Primary limitation}: leakage into $n\geq 2$ Fock states driven
        by the Displacement gates.  Smaller $|\alpha|$ reduces leakage at the
        cost of less Fock-mixing power.
  \item \textbf{Decoherence}: the 7.1-\si{\micro s} total gate time is
        comparable to the qubit $T_1\approx9.8$~\si{\micro s}, making coherence
        the dominant practical constraint.
\end{itemize}

\subsection{Physical Insight: Fock-Mixing Constraint}

A central finding is the \emph{Fock-mixing constraint}: SQR gates (qubit
rotations conditioned on cavity Fock number) can rotate the qubit within each
Fock sector but cannot transfer population between Fock levels.  Since
$\Utarget$ contains $H_c$ and $\mathrm{CNOT}_{q\to c}$, both of which mix
the $\ket{n=0}$ and $\ket{n=1}$ sectors, Displacement is a necessary primitive
in any synthesis library.  This is confirmed by the SQR+CP maximum of $F=0.5$
(Phase~2) and by the success of libraries including Displacement.

\subsection{Comparison of Libraries}

Library 1 (D+$R_q$+SQR) outperforms Library 2 (D+$R_q$+CP) at physical
parameters.  The SQR gate provides full Fock-conditioned qubit rotations
on the Bloch sphere, whereas the CP gate applies only Fock-selective phases
with no qubit rotation, giving it fewer degrees of freedom for synthesizing
the complex entangling structure of $\Utarget$.  Both libraries require
Displacement for Fock-level mixing, confirming the Fock-mixing constraint.

\subsection{Robustness}

The L1c (D+$R_q$+SQR) sequence is robust to the most common experimental
calibration errors: zero fidelity sensitivity to $\chi$ miscalibration
and pulse duration errors (both entering as $\chi t$, absorbed by variational
phase optimization), and small sensitivity to amplitude and phase errors.
This robustness arises because the dispersive drift contribution to the
logical subspace ($n=0$ and $n=1$) is fully absorbed into the variational
SQR azimuthal angles $\phi_n$.

%============================================================
\appendix
\section{Code and Data}
%============================================================

All scripts are located in \texttt{scripts/}, data in \texttt{data/}, and
figures in \texttt{figures/} relative to the study root directory.

\begin{table}[H]
  \centering
  \caption{Scripts and data files.}
  \begin{tabular}{ll}
    \toprule
    File & Content \\
    \midrule
    \texttt{scripts/phase1\_algebraic.py} & Algebraic checks, CNOT factorization \\
    \texttt{scripts/phase2\_idealized.py} & 4D ideal synthesis, depth sweep \\
    \texttt{scripts/phase3\_embedding.py} & Truncation convergence, leakage \\
    \texttt{scripts/phase4\_variational.py} & Main variational synthesis \\
    \texttt{scripts/phase5\_pulses.py} & Pulse effects, coherence budget \\
    \texttt{scripts/phase6\_robustness.py} & Robustness sweeps \\
    \midrule
    \texttt{data/phase\{1..6\}\_results.json} & Numerical results \\
    \texttt{figures/phase\{1..6\}\_*.pdf} & Publication-quality figures \\
    \bottomrule
  \end{tabular}
\end{table}

\end{document}
